%% LyX 2.5.1 created this file.  For more info, see https://www.lyx.org/.
%% Do not edit unless you really know what you are doing.
\documentclass[english]{article}
\usepackage{courier}
\usepackage[T1]{fontenc}
\usepackage[latin9]{inputenc}
\usepackage{babel}
\usepackage{float}
\usepackage{enumitem}
\usepackage{amsmath}
\usepackage{amsthm}
\usepackage{graphicx}
\usepackage{geometry}
\geometry{verbose,tmargin=1in,bmargin=1in,lmargin=1in,rmargin=1in}
\usepackage[pdfusetitle,
 bookmarks=true,bookmarksnumbered=false,bookmarksopen=false,
 breaklinks=false,pdfborder={0 0 1},backref=false,colorlinks=false]
 {hyperref}

\makeatletter

%%%%%%%%%%%%%%%%%%%%%%%%%%%%%% LyX specific LaTeX commands.
%% A simple dot to overcome graphicx limitations
\newcommand{\lyxdot}{.}


%%%%%%%%%%%%%%%%%%%%%%%%%%%%%% Textclass specific LaTeX commands.
\newlength{\lyxlabelwidth}      % auxiliary length 

\makeatother

\theoremstyle{plain}
\newtheorem{thm}{\protect\theoremname}
\newtheorem{question}[thm]{\protect\questionname}
\theoremstyle{definition}
\newtheorem*{sol*}{\protect\solutionname}
\usepackage{listings}
\providecommand{\questionname}{Question}
\providecommand{\solutionname}{Solution}
\providecommand{\theoremname}{Theorem}
\renewcommand{\lstlistingname}{Listing}

\begin{document}
\title{Assignment 2 - Policy Gradient vs Value-Based RL in JAX}
\author{Jonathan Horton (\#7710257)}
\maketitle

\section{Note}

In general in this report, the results discussed are for the case
where the grid world is empty and no obstacles are present, unless
otherwise specified. Only a goal and a starting point exist.

\section{Computer/GPU Details }

The code was run on a computer with 32 GB of RAM, a AMD Ryzen 9 3900X
12-Core 3.80GHz Processor, and on a NVIDIA GeForce RTX 2080 (8GB)
GPU.

\section{Description of Environment}

Gymnasium was used for the environment due to it being a more widely
used library and having a wider range of tutorials and documentation
available online, thereby simplifying implementation. The environment
is a simple 7 by 7 Grid. The agent starts at the top right corner
at the point $\left(1,1\right)$ of the environment and works towards
the point $\ensuremath{\left(5,5\right)}$. It will not move if it
attempts to move off of any of the edges of the grid. The agent cannot
move diagonally; it can only move (i.e. select the action) either
up, down, left, or right. The experiment was run twice, once with
pits and once without. A pit adds a negative reward when the agent
traverses onto that square and resets the position of the agent to
the point $\ensuremath{\left(1,1\right)}$ when entered. The episode
continues even if a pit has been entered.

Observations provided to the reinforcement algorithms are: the normalized
coordinates, the Manhattan distance to the goal, and the direction
to the goal.

A small reward is given if a step moves the agent closer to the goal,
and a small penalty is given for moving away. In addition, another
constant small penalty is added to the reward encourage shorter paths.
A large penalty is added to the reward if the agent steps in a pit.
Lastly, A large reward is provided for reaching the end goal.

A wrapper using \lstinline[language=Python,basicstyle={\ttfamily}]!gym.Wrapper!
was implemented

The research questions that are being investigated in this assignment:
\begin{enumerate}
\item Make a comparison of DQN vs REINFORCE-wb in a blank environment. As
required by the assignment, 10 different seeds are tried for each
algorithm and set of hyperparameter settings (Learning rates: $\ensuremath{0.001}$,
$\ensuremath{0.01}$, and $\ensuremath{0.1}$; Discount factors: $\ensuremath{\gamma=0.9}$,
$\ensuremath{0.99}$).
\item Observe the effects of a vertical pit (as defined above) on the learning
rates of DQN and REINFORCE-wb.
\end{enumerate}

\section{Learning Curves for All Agents}

All agents were trained for 200 episodes at a maximum of 100 steps.
This report will be submitted with all of the learning curves for
all of the agents. Note that they will not all be shown in this report
as that would take an unreasonable amount of space in the report,
as there are 48 graphs. Some examples of learning curves are shown
in Figure \ref{fig:REINFORCE-wb:-Reward-vs} and Figure \ref{fig:DQN:-Reward-vs}.
From these two graphs, it can be observed that REINFORCE is much less
noisy than DQN, though both suffer from spikes in rewards during the
training process. Generally speaking, with the specific hyperparameters
shown here, both are able to learn a rewarding behaviour for the defined
grid world environment. Note that there are 10 subplots here in each
graph: one for each of the seeds used in testing. The random baseline
is shown in red in all graphs for each seed. The overlayed version
of these graphs are shown in Figure \ref{fig:REINFORCE-wb-overlayed:-Reward}
and Figure \ref{fig:DQN-overlayed:-Reward}. The increased noisy behavious
of DQN is particularly visible in Figure \ref{fig:DQN-overlayed:-Reward}.

\begin{figure}[H]
\begin{centering}
\includegraphics[scale=0.4]{results/reinforce_gamma0\lyxdot 99_lr0\lyxdot 001_reward_subplots}
\par\end{centering}
\caption{\label{fig:REINFORCE-wb:-Reward-vs}REINFORCE-wb: Reward vs Episodes
($\ensuremath{\gamma}$=0.99, lr=0.001)}

\end{figure}

\begin{figure}[H]
\begin{centering}
\includegraphics[scale=0.4]{results/dqn_gamma0\lyxdot 99_lr0\lyxdot 001_reward_subplots}
\par\end{centering}
\caption{\label{fig:DQN:-Reward-vs}DQN: Reward vs Episodes ($\ensuremath{\gamma}$=0.99,
lr=0.001)}
\end{figure}

\begin{figure}[H]
\begin{centering}
\includegraphics[scale=0.4]{results/reinforce_gamma0\lyxdot 99_lr0\lyxdot 001_reward_overlay}
\par\end{centering}
\caption{\label{fig:REINFORCE-wb-overlayed:-Reward}REINFORCE-wb overlayed:
Reward vs Episodes ($\ensuremath{\gamma}$=0.99, lr=0.001)}
\end{figure}

\begin{figure}[H]
\begin{centering}
\includegraphics[scale=0.4]{results/dqn_gamma0\lyxdot 99_lr0\lyxdot 001_reward_overlay}
\par\end{centering}
\caption{\label{fig:DQN-overlayed:-Reward}DQN overlayed: Reward vs Episodes
($\ensuremath{\gamma}$=0.99, lr=0.001)}
\end{figure}


\section{Policy Visualization}

All agents were trained for 200 episodes at a maximum of 100 steps.
Some examples of policy visualization are shown in Figures \ref{fig:DQN-Final-Policy:},
\ref{fig:REINFORCE-wb:-Final-policy}, \ref{fig:DQN-Final-policy:}
and \ref{fig:REINFORCE-wb:-Final-policy-1}. Figures \ref{fig:DQN-Final-Policy:}
and \ref{fig:REINFORCE-wb:-Final-policy} show the policies as arrows,
whereas Figures \ref{fig:DQN-Final-policy:} and \ref{fig:REINFORCE-wb:-Final-policy-1}
use colours. Regardless, both show the same information: the final
policy indicating the optimal action to take in each state. The agent
starts at the point labeled with the blue S and ends at the point
labeled with a red G. The colours shown in Figure \ref{fig:DQN-Final-policy:}
and \ref{fig:REINFORCE-wb:-Final-policy-1} correspond to the following:
yellow for left, blue for right, purple for up, and greed for down.

Some of the illustrated policies would always give the solution regardless
of the starting action taken and assuming a greedy policy. For example,
Seed 5 in figure \ref{fig:DQN-Final-Policy:} is optimal. Others do
not meet this criteria; for example, seed 4 in Figure \ref{fig:DQN-Final-Policy:}
heads eventually left into a wall after going down at the starting
point.

\begin{figure}[H]
\begin{centering}
\includegraphics[scale=0.4]{results/dqn_gamma0\lyxdot 99_lr0\lyxdot 001_policy_arrows}
\par\end{centering}
\caption{\label{fig:DQN-Final-Policy:}DQN Final Policy: arrows($\ensuremath{\gamma}$=0.99,
lr=0.001)}
\end{figure}

\begin{figure}[H]
\begin{centering}
\includegraphics[scale=0.4]{results/reinforce_gamma0\lyxdot 99_lr0\lyxdot 001_policy_arrows}
\par\end{centering}
\caption{\label{fig:REINFORCE-wb:-Final-policy}REINFORCE-wb: Final policy
arrows($\ensuremath{\gamma}$=0.99, lr=0.001)}
\end{figure}

\begin{figure}[H]
\begin{centering}
\includegraphics[scale=0.4]{results/dqn_gamma0\lyxdot 99_lr0\lyxdot 001_policy_heatmaps}
\par\end{centering}
\caption{\label{fig:DQN-Final-policy:}DQN Final policy: Heatmap ($\ensuremath{\gamma}$=0.99,
lr=0.001)}
\end{figure}

\begin{figure}[H]
\begin{centering}
\includegraphics[scale=0.4]{results/reinforce_gamma0\lyxdot 99_lr0\lyxdot 001_policy_heatmaps}
\par\end{centering}
\caption{\label{fig:REINFORCE-wb:-Final-policy-1}REINFORCE-wb: Final policy
heat maps ($\ensuremath{\gamma}$=0.99, lr=0.001)}
\end{figure}


\section{Hyperparameters Sensitivity Analysis}

Overall, both algorithms are sensitive to hyperparameters. It is possible
that learning would occur if the algorithms were run for a greater
number of episodes; this was not explored in this assignment. 

DQN learned very poorly when the learning rate was set to $\ensuremath{0.1}$
when compared to its performance at $\ensuremath{0.01}$ and $\ensuremath{0.001}$
for gamma = 0.9. Moreover, for DQN, increasing the gamma value to
$\ensuremath{0.99}$ dramatically reduced learning across all learning
rates, except for 0.001 where the results were similar but noisier.
In general, a smaller value of gamma and a smaller value for the learning
rate produces better results.

When the gamma value was set to 0.9, REINFORCE-wb performed quite
well at a small learning rate of 0.001, but as the value increased,
it fails to learn. A very similar behaviour is also observed when
the gamma value is set to 0.99, suggesting that REINFORCE-wb is less
sensitive to this hyperparameter than DQN. In general, smaller values
for both hyperparameters were also preferred for REINFORCE-wb.

\section{Reflections}
\begin{question}
The mean performance of your agent versus the variability across seeds.
If there is a high-degree of variability discuss the potential causes
and how you could assess this.
\begin{sol*}
For the most part, there was not a large amount of variability across
seeds. However, a particular exception to this rule was REINFORCE-WB
at $\ensuremath{\gamma=0.99}$ and $\ensuremath{\text{lr}=0.01}$
and DQN at $\ensuremath{\gamma=0.9}$ and $\ensuremath{\text{lr}=0.1}$.
In the case of REINFORCE-wb at t $\ensuremath{\gamma=0.99}$ and $\ensuremath{\text{lr}=0.01}$
a possible cause might be that the long-horizon credit assignment
combined with the high learning rate makes return estimates high-variance.
Thus, tiny differences in early trajectories across seeds could result
in different policy updates. In the case of DQN, the large learning
rate might be causing Q-updates to overshoot, which causes value estimates
to be higher in variance, thereby having the entire process have a
higher sensitivity to the exact replay samples each seed encounters.
\end{sol*}
\end{question}

\begin{question}
How JAX\textquoteright s functional style influenced your implementation.
\begin{sol*}
JAX's function style impacted how both my training loops and environment
was structured. Some modifications to the environment were necessary
to ensure optimal performance. Moreover, because functions must remain
pure, parameters, states, and random seeds were seperated to avoid
having hidden mutable states. Loss functions were particularly impacted
by JAX because they accept parameters as inputs and return losses
and other values, making the update step a transformation rather than
an in-place mutation (as other paradigms use). The functionnal constraints
also made it easy to implement both vmap and JIT. This changed how
batches had to be organized as well as policy evaluation, so that
both could be expressed as pure array transformations. 
\end{sol*}
\end{question}

\begin{question}
Observe the effects of a vertical pit (as defined above) on the learning
rates of DQN and REINFORCE-wb.
\begin{sol*}
Frankly, the addition of the pit seems to have had minimal effects
on training. The pit was placed at (0,2), (1,2), (2,2), (3,2), and
(4,2). Performance still drops as gamma increases and lr increases.
The best results, shown are in Figures \ref{fig:REINFORCE-wb:-Reward-vs-1}
and \ref{fig:DQN:-Reward-vs-1}. The random agent performs poorly
because of the negative rewards provided by the pit. The arrow diagrams
suggest that an error may have occured during saving, or that the
agent is learning a loop to obtain rewards for this particular case.
\end{sol*}
\begin{figure}[H]
\begin{centering}
\includegraphics[scale=0.4]{resultsWithPits/dqn_gamma0\lyxdot 99_lr0\lyxdot 001_reward_subplots}
\par\end{centering}
\caption{\label{fig:REINFORCE-wb:-Reward-vs-1}DQN: Reward vs Episodes ($\ensuremath{\gamma}$=0.99,
lr=0.001)}
\end{figure}

\begin{figure}[H]
\begin{centering}
\includegraphics[scale=0.4]{resultsWithPits/reinforce_gamma0\lyxdot 99_lr0\lyxdot 001_reward_subplots}
\par\end{centering}
\caption{\label{fig:DQN:-Reward-vs-1}REINFORCE-wb: Reward vs Episodes ($\ensuremath{\gamma}$=0.99,
lr=0.001)}
\end{figure}

\end{question}


\section{AI Disclosure and Note}

The Syllabus states ``Assignments \& Presentations: Use of AI tools
(e.g., ChatGPT, Gemini, Copilot) are not permitted.'' As such, my
response to the questions ``What did AI suggest? What did you accept/reject
and why?'' is that AI use is not permitted for assignments (per the
syllabus), and that therefore it was not used. 

If the policy for the course regarding use of AI in assignments has
changed, I would request that you please make an announcement clarifying
in what ways AI can be used for assignments.
\end{document}
